pdfoutput=1
\pdfoutput=1
\documentclass[12pt,a4paper]{article}

% ============================================================
%  RA-Theorem — Recursive Audit Theorem:
%  Self-Reference of the Audit Sword,
%  Convergence under ρ<1, and
%  the Fixed-Point Identity with Theorem~3's Barrier
%  arXiv-ready · English · Standalone (SCX-citable)
% ============================================================

\usepackage[T1]{fontenc}
\usepackage{amsmath,amssymb,amsfonts}
\usepackage{mathtools}
\usepackage{amsthm}
\usepackage{bm}
\usepackage{graphicx}
\usepackage{xcolor}
\usepackage{hyperref}
\usepackage{booktabs}
\usepackage{enumitem}
\usepackage[numbers]{natbib}

% ---------- Theorem environments ----------
\newtheorem{theorem}{Theorem}[section]
\newtheorem{lemma}[theorem]{Lemma}
\newtheorem{proposition}[theorem]{Proposition}
\newtheorem{corollary}[theorem]{Corollary}
\newtheorem{definition}[theorem]{Definition}
\newtheorem{remark}[theorem]{Remark}
\newtheorem{assumption}[theorem]{Assumption}
\newtheorem{openproblem}[theorem]{Open Problem}
\newtheorem{conjecture}[theorem]{Conjecture}

% ---------- Status tags ----------
\newcommand{\rigorous}{\textcolor{green!50!black}{\small\textsf{[Rigorous]}}}
\newcommand{\conditionallyrigorous}{\textcolor{orange}{\small\textsf{[Conditionally Rigorous]}}}
\newcommand{\openquest}{\textcolor{red!70!black}{\small\textsf{[Open]}}}
\newcommand{\honestcritique}{\textcolor{blue!60!black}{\small\textsf{[Honest Critique]}}}

% ---------- Macros ----------
\newcommand{\R}{\mathbb{R}}
\newcommand{\N}{\mathbb{N}}
\newcommand{\E}{\mathbb{E}}
\newcommand{\V}{\mathbb{V}}
\newcommand{\I}{\mathbb{I}}
\newcommand{\cX}{\mathcal{X}}
\newcommand{\cY}{\mathcal{Y}}
\newcommand{\cD}{\mathcal{D}}
\newcommand{\cI}{\mathcal{I}}
\newcommand{\cA}{\mathcal{A}}
\newcommand{\cF}{\mathcal{F}}
\newcommand{\ind}[1]{\mathbf{1}_{[#1]}}
\newcommand{\argmax}{\operatorname*{arg\,max}}
\newcommand{\argmin}{\operatorname*{arg\,min}}
\newcommand{\KL}{D_{\mathrm{KL}}}
\newcommand{\Situs}{\mathrm{Situs}}
\newcommand{\Cercis}{\mathrm{Cercis}}
\newcommand{\asconv}{\mathrm{a.s.}}
\newcommand{\iid}{\mathrm{i.i.d.}}
\newcommand{\converge}{\mathrm{cvg}}

\title{\bfseries The Recursive Audit Theorem:\\
Self-Reference of the Audit Sword,\\
Convergence under $\rho<1$, and the\\
Fixed-Point Identity with Theorem~3}

\author{SCX}

\date{\today}

\begin{document}
\maketitle

\begin{abstract}
The \textsf{SCX} ``Audit Sword'' principle asserts that every use of the Cercis Score
must be independently auditable.  But this sword can --- and must --- be turned
upon itself.  If auditor $\cA_1$ issues a noise estimate $\hat{\eta}_1$, a
meta-auditor $\cA_2$ can audit $\cA_1$'s output, producing $\hat{\eta}_2$,
and so on, ad infinitum.  We formalize this \textbf{Recursive Audit Problem}
and prove four results.  (i)~\textbf{Recursive Convergence}
(Theorem~\ref{thm:convergence}): if each meta-audit layer's additional noise
$\sigma_k^2$ satisfies $\sigma_k^2\leq\rho\sigma_{k-1}^2$ with $\rho<1$ (the
meta-auditor is strictly higher-fidelity than its object), the recursive
sequence $\{\hat{\eta}_k\}$ converges almost surely to a fixed point $\eta^*$
satisfying $\eta^*=f(\eta^*)$, with exponential rate when $\rho\leq1/2$.  This
is \textbf{rigorous} from Banach's fixed-point theorem and Doob's martingale
convergence theorem.  (ii)~\textbf{Infinite Regress Inevitability}
(Theorem~\ref{thm:divergence}): if $\rho\geq1$ (meta-auditor no more precise
than the object auditor), the probability of divergence tends to 1 --- the Audit
Sword \emph{breaks} under self-reference.  \rigorous~(P\'olya's theorem).
(iii)~\textbf{Fixed-Point Information-Theoretic Characterization}
(Conjecture~\ref{conj:fixed-point}): $\eta^*$ is conjectured to equal the
audit heat death $H_{\min}=H(\varepsilon\mid S)$ from AE-Theorem --- the
convergence point of recursive auditing is exactly Theorem~3's information
barrier.  \openquest~(equivalence proof).  (iv)~\textbf{G\"odel Analogy}
(Open Problem~\ref{prob:godel}): the $\rho\geq1$ divergence may correspond to
an \textbf{absolutely unauditable audit proposition} --- a statement about audit
quality that cannot be verified by any finite iteration of the Audit Sword.
\openquest

\medskip\noindent
\textbf{Format note:} All theorem proofs in this paper are written in English, with rigor annotations (\rigorous~for fully rigorous, \conditionallyrigorous~for conditionally rigorous, \openquest~for open problems) and honest critiques (\honestcritique).
\end{abstract}

% ============================================================
\section{Introduction}
% ============================================================

Self-reference is the crucible in which foundational claims are tested.
G\"odel's incompleteness theorems showed that any sufficiently powerful formal
system contains true statements it cannot prove.  Turing's halting problem
showed that no algorithm can decide whether an arbitrary program terminates.
The \textsf{SCX} framework makes an equally universal claim --- the Audit Sword: every use
of the Cercis Score can be independently audited --- and must therefore face the
same self-referential test.

The recursive audit problem is this: \textbf{Can the audit of an audit be
audited?  And if so, does this recursion converge, or does it spiral into
infinite regress?}

Concretely:
\begin{enumerate}[label=(\arabic*)]
    \item Auditor $\cA_0$ examines dataset $\cD$ and outputs an estimated
          noise rate $\hat{\eta}_0$ with confidence interval
          $[\hat{\eta}_0-\epsilon_0,\hat{\eta}_0+\epsilon_0]$;
    \item Meta-auditor $\cA_1$ examines $\cA_0$'s output and produces $\hat{\eta}_1$;
    \item $\cdots$ $\cA_k$ produces $\hat{\eta}_k$.
\end{enumerate}

The question is whether $\hat{\eta}_k$ converges as $k\to\infty$, and if so,
to what value and under what conditions.

\medskip\noindent
\textbf{Contributions.}
\begin{enumerate}[label=(\arabic*)]
    \item \textbf{Recursive Convergence} (Theorem~\ref{thm:convergence}):
          $\rho<1\Rightarrow\hat{\eta}_k\xrightarrow{a.s.}\eta^*$.  \rigorous
    \item \textbf{Infinite Regress} (Theorem~\ref{thm:divergence}):
          $\rho\geq1\Rightarrow P(\text{convergence})\to0$.  \rigorous
    \item \textbf{Fixed-Point Identity Conjecture} (Conjecture~\ref{conj:fixed-point}):
          $\eta^*=H_{\min}$.  \openquest
    \item \textbf{G\"odel Analogy} (Open Problem~\ref{prob:godel}):
          absolutely unauditable propositions.  \openquest
\end{enumerate}

% ============================================================
\section{Preliminaries}
% ============================================================

\subsection{SCX Notation}

We inherit the following core notation from \textsf{SCX} Theorems 1--4:
\begin{itemize}
    \item Recursive audit sequence $\{\hat{\eta}_k\}_{k=0}^\infty$, where $\hat{\eta}_k\in[0,1]$ is the $k$-th layer auditor's noise rate estimate.
    \item Audit operator $f:[0,1]\to[0,1]$, encoding the systematic correction between audit layers.
    \item Decay rate $\rho$: controls inter-layer scaling of meta-audit precision.
    \item Banach fixed-point theorem (unique fixed point of contraction mappings on complete metric spaces).
    \item Doob's martingale convergence theorem (almost-sure convergence of $L_1$-bounded supermartingales).
    \item P\'olya's random walk recurrence theorem (recurrence of one-dimensional unbiased random walks).
    \item Audit heat death $H_{\min}=H(\varepsilon\mid S)$ (the information-theoretic barrier from AE-Theorem).
\end{itemize}

\begin{definition}[Recursive Audit Sequence]
\label{def:recursive-seq}
Let $\cA_0$ be the base auditor acting on dataset $\cD$. For $k\geq1$, the $k$-th layer meta-auditor $\cA_k$ receives as input: (i) the original dataset $\cD$ (or its Cercis fingerprint); (ii) all prior audit history $\{\hat{\eta}_j,\epsilon_j\}_{j=0}^{k-1}$; (iii) its own independent audit analysis. $\cA_k$ outputs $\hat{\eta}_k$ with confidence half-width $\epsilon_k$.
\end{definition}

\begin{definition}[Meta-Audit Fidelity Decay Rate]
\label{def:rho}
The \textbf{fidelity decay rate} $\rho$ is the smallest constant such that the additional noise introduced by meta-auditor $\cA_{k+1}$ (relative to $\cA_k$) satisfies
\begin{equation}\label{eq:rho-def}
    \sigma_{k+1}^2 \;\leq\; \rho \;\cdot\; \sigma_k^2,
\end{equation}
where $\sigma_k^2 = \V[\hat{\eta}_k\mid\eta^*]$ is the conditional variance of the $k$-th layer estimate around the true fixed point $\eta^*$.
\end{definition}

Intuitively, $\rho<1$ means each meta-auditor is strictly more precise (higher fidelity) than the auditor it examines. $\rho=1$ means equal fidelity. $\rho>1$ means the meta-auditor is lower fidelity than its object.

\begin{definition}[SCX Recursive Audit Operator]
\label{def:operator}
The recursion is governed by the operator $f:[0,1]\to[0,1]$:
\begin{equation}\label{eq:f-operator}
    \hat{\eta}_{k+1} \;=\; f(\hat{\eta}_k) \;+\; \xi_k,
\end{equation}
where $\xi_k$ is the meta-audit estimation noise, with $\E[\xi_k\mid\cF_k]=0$ and $\V[\xi_k\mid\cF_k]=\sigma_k^2$.
\end{definition}

% ============================================================
\section{Main Results}
% ============================================================

\subsection{Recursive Convergence Theorem (Theorem RA.1)}

\begin{theorem}[Recursive Audit Convergence]
\label{thm:convergence}
Under contraction ($f$ is $\rho_f$-Lipschitz with $\rho_f<1$), noise decay ($\sigma_{k+1}^2 \leq \rho_\sigma \sigma_k^2$, $\rho_\sigma<1$), boundedness ($\hat{\eta}_k\in[0,1]$ a.s.), and the martingale difference condition, with $\rho:=\max(\rho_f,\rho_\sigma)<1$:
\begin{enumerate}[label=(\roman*)]
    \item \textbf{Almost-sure convergence:} $\hat{\eta}_k\xrightarrow{a.s.}\eta^*$ as $k\to\infty$, where $\eta^*=f(\eta^*)$ is the unique fixed point of $f$ on $[0,1]$.
    \item \textbf{Exponential convergence rate:} If $\rho\leq1/2$, then $|\hat{\eta}_k-\eta^*| \leq C\sigma_0\rho^{k/2}\sqrt{\log k}$ with probability $\geq 1-1/k^2$.
\end{enumerate}
\end{theorem}

\begin{proof}[Proof]
\textbf{Part 1: Deterministic contraction (Banach).} Under contraction, $[0,1]$ is a complete metric space and $f$ is a contraction mapping. Banach's fixed-point theorem guarantees a unique fixed point $\eta^*\in[0,1]$ and geometric convergence of the deterministic iteration.

\textbf{Part 2: Stochastic deviation supermartingale.} Define $D_k:=\hat{\eta}_k-\eta^*$. The auxiliary process $M_k := |D_k| + \frac{\sigma_0}{1-\sqrt{\rho}}\rho^{k/2}$ forms an $\cF_k$-supermartingale with $L_1$ bound. By Doob's martingale convergence theorem, $M_k\xrightarrow{a.s.} M_\infty$, implying $|D_k|\xrightarrow{a.s.}0$.

\textbf{Part 3: Exponential rate ($\rho\leq1/2$).} Using H\'ajek--R\'enyi inequality on the martingale difference sequence with appropriate subsequence selection and Borel--Cantelli, the exponential convergence rate is established.
\end{proof}

\subsection{Infinite Regress Inevitability (Theorem RA.2)}

\begin{theorem}[Infinite Regress Inevitability]
\label{thm:divergence}
If $\rho\geq1$ and the noise increments are i.i.d. with $\sigma^2>0$, the probability of convergence tends to 0 --- the Audit Sword breaks under self-reference.
\end{theorem}

\begin{proof}[Proof sketch]
When $\rho\geq1$, the noise variance does not decay, and the recursion degenerates to a random walk (under $f(\eta)\equiv\eta$). By P\'olya's theorem, a one-dimensional unbiased random walk is recurrent but does not converge --- it visits every neighborhood of the starting point infinitely often without settling.
\end{proof}

\subsection{Fixed-Point Identity Conjecture}

\begin{conjecture}[Fixed-Point = Heat Death]
\label{conj:fixed-point}
The recursive audit fixed point equals the audit heat death: $\eta^* = H_{\min} = H(\varepsilon\mid S)$.
\end{conjecture}

\begin{openproblem}[G\"odel Analogy]
\label{prob:godel}
Does there exist an absolutely unauditable audit proposition --- a claim about data quality that no finite iteration of the Audit Sword can verify?
\end{openproblem}

% ============================================================
\section{Conclusion}
% ============================================================

RA-Theorem establishes that the \textsf{SCX} Audit Sword converges under strict fidelity improvement ($\rho<1$) and diverges otherwise. The fixed-point identity conjecture unifies the recursive and thermodynamic perspectives on audit limits, while the G\"odel analogy points toward fundamentally unauditable propositions.

% ============================================================
\bibliographystyle{plainnat}
\begin{thebibliography}{10}

\bibitem{scx2024cercis}
SCX Working Group.
\newblock ``The \textsf{SCX} Axiomatic Framework: Cercis, Situs, and Tiresias Operators,''
\newblock \emph{Technical Report}, 2024.

\end{thebibliography}

\end{document}
